\color{unidarkgreen}{\subsection{Контроль силового выключателя (КСВ)}}
\color{black}

\begin{enumerate}[label=\arabic{section}.\arabic{subsection}.\arabic{enumi}, labelsep=4pt, leftmargin=0pt, itemindent=57pt, itemsep=0pt, parsep=0pt] % если что можно крутить itemsep

\item
Функциональный блок контроля силового выключателя (КСВ) предназначен для обработки технологической
информации от выключателя и его цепей управления, а также от трансформатора тока, позволяя определять их состояние,
информировать о возможных неисправностях и формировать соответствующую сигнализацию.

\item
Контроль целостности цепей управления осуществляется по сигналам состояния цепей первой и второй
группы электромагнитов отключения (ЭМО1 и ЭМО2) и цепи электромагнита включения (ЭМВ). Наличие второй
группы электромагнитов отключения определяется положением программного переключателя \textbf{«Кол-во ЭМО
Вn»}. При одновременном наличии или отсутствии сигналов контроля цепей ЭМО и ЭМВ в течение
регулируемой выдержки времени \textbf{«Тср Неиспр. ЦУ Вn»} формируется сигнал \textit{«Неиспр. ЦУ Вn»}. Выдержка
времени должна быть не менее времени взвода пружины привода выключателя.

\item
Контроль готовности привода выключателя осуществляется по внешнему сигналу \textit{«Пруж.
не заведены Вn»}. Сигнал о незаведенном состоянии пружины \textit{«Ср. Пруж. не заведены Вn»} формируется с выдержкой
времени \textbf{«Тср Пруж. не заведены Вn»}. 

\item
Предусмотрен контроль давления элегаза в баке выключателя по внешним сигналам \textit{«Низк. давл. SF6
Вn»} и \textit{«Авар. давл. SF6 Вn»} от специальных датчиков давления. Соответствующие сигналы о низком и
аварийном давлении элегаза \textit{«Ср. Низк. давл. SF6 Вn»} и \textit{«Ср. Авар. давл. SF6 Вn»} формируются с выдержками
времени, равными 1 с.

\item
Контроль положения автомата шинок питания осуществляется по внешнему сигналу \textit{«Автомат ШП Вn»}.
Сигнал об отключенном положении автомата \textit{«Ср. Автомат ШП Вn»} формируется с выдержкой времени выдержкой времени,
равной 1 с. 

\item
Обогрев выключателя контролируется по внешнему сигналу \textit{«Неиспр. Обогр. Вn»}. Соответствующий
сигнал о неисправности обогрева \textit{«Ср. Неиспр. Обогр. Вn»} формируется с регулируемой выдержкой времени
\textbf{«Тср Неиспр. Обогр. Вn»}.

\item
Сигнализация о неисправности выключателя (сигнал \textit{«Неиспр. Вn»}) срабатывает при наличии одного из
сигналов: \textit{«Неиспр. ЦУ Вn»}, \textit{«Ср. Пруж. не заведены Вn»}, \textit{«Ср. Низк. давл. SF6 Вn»}, \textit{«Ср. Авар. давл. SF6 Вn»},
\textit{«Ср. Автомат ШП Вn»}, \textit{«Ср. Неиспр. обогр. Вn»}, а также от сигнала о непереключении фаз \textit{«Ср. ЗНФ Вn»}.

\vspace{2mm}
\begin{figure}[H]
\centering
\includegraphics[width=1\textwidth,height=1\textheight,keepaspectratio]{\fbpath/761.3371 КСВ/img/CTRL_Q.pdf}
\captionof{figure}{Функциональный блок контроля исправности выключателя В1}\label{pict_ksv:CTRL_Q}
\end{figure}

\item
В устройстве также предусмотрен контроль состояния ТТ. Контроль давления элегаза в ТТ
осуществляется по внешним сигналам \textit{«Низк. давл. SF6 ТТ Вn»} и \textit{«Авар. давл. SF6 ТТ Вn»}. Соответствующие
сигналы о низком и аварийном давлении элегаза \textit{«Ср. Низк. давл. SF6 Вn»} и \textit{«Ср. Авар. давл. SF6 Вn»}
формируются с выдержками времени, равными 1 с. Обогрев ТТ контролируется по внешнему сигналу \textit{«Неиспр.
обогр. ТТ В1»}. Соответствующий сигнал о неисправности обогрева \textit{«Ср. Неиспр. обогр. ТТ Вn»} формируется
с регулируемой выдержкой времени \textbf{«Тср Неиспр. обогр. ТТ Вn»}. По наличию одного из указанных сигналов
формируется сигнал о неисправности ТТ \textit{«Неиспр. ТТ Вn»}. При вводе программного ключа \textbf{«Откл. от Авар.
давл. SF6 ТТ Вn»} предусматривается действие на отключение выключателя при аварийном снижении
давления элегаза в ТТ.

\vspace{2mm}
\begin{figure}[H]
\centering
\includegraphics[width=1\textwidth,height=1\textheight,keepaspectratio]{\fbpath/761.3371 КСВ/img/CTRL_CT_Q.pdf}
\captionof{figure}{Функциональный блок контроля исправности ТТ В1}\label{pict_ksv:CTRL_CT_Q}
\end{figure}

\item
Контроль отсутствия питания цепей сигнализации выключателя и ТТ осуществляется по внешним
сигналам \textit{«Контр. опертока ЭМВ/ЭМО1 Вn»}, \textit{«Контр. опертока ЭМО2 Вn»} и \textit{«Контр. опертока сигн. Вn»}. При
состоянии «0» одного из сигналов срабатывает сигнализация о неисправности цепей сигнализации
выключателя и ТТ (сигнал \textit{«Неиспр. сигн. Вn/ТТ Вn»}) с выдержкой времени 1 с. 

\vspace{2mm}
\begin{figure}[H]
\centering
\includegraphics[width=1\textwidth,height=1\textheight,keepaspectratio]{\fbpath/761.3371 КСВ/img/CTRL_OC_Q.pdf}
\captionof{figure}{Функциональный блок контроля оперативного тока цепей сигнализации В1}\label{pict_ksv:CTRL_OC_Q}
\end{figure}

\item
Для электромагнитов управления предусмотрена защита от длительного протекания тока. Контроль
протекания тока через ЭМВ, ЭМО1 и ЭМО2 осуществляется при помощи токовых реле, положения контактов
которых должны быть заведены на соответствующие входные сигналы ФСУ «Работа ЭМВ Вn», «Работа ЭМО1
Вn», «Работа ЭМО2 Вn». При длительном протекании тока по цепям ЭМО1 и ЭМВ формируется сигнал «Обест.
ЭМО1 и ЭМВ Вn», действующий на независимый расцепитель автомата питания цепей ЭМО1 и ЭМВ.

\item
Аналогично при длительном протекании тока по цепи ЭМО2 формируется сигнал \textit{«Обест. ЭМО2 Вn»} для
действия на независимый расцепитель автомата питания цепи ЭМО2. Регулируемые выдержки времени
\textbf{«Тср защиты ЭМВ Вn»}, \textbf{«Тср защиты ЭМО1 Вn»} и \textbf{«Тср защиты ЭМО2 Вn»} определяют время
срабатывания на обесточивание ЭМВ, ЭМО1 и ЭМО2 соответственно. Кроме того, наличие сигнала \textit{«Ср. Авар.
давл. SF6 Вn»} действует на обесточивание всех электромагнитов.

\vspace{2mm}
\begin{figure}[H]
\centering
\includegraphics[width=0.9\textwidth,height=0.9\textheight,keepaspectratio]{\fbpath/761.3371 КСВ/img/PROT_EMU_Q.pdf}
\captionof{figure}{Функциональный блок защиты ЭМУ В1 от длительного протекания тока}\label{pict_ksv:PROT_EMU_Q}
\end{figure}

\item
На блокировку отключения выключателя действует сигнал \textit{«Ср. Авар. давл. SF6 Вn»}.

\item
На блокировку включения действуют сигналы \textit{«Ср. Авар. давл. SF6 Вn»}, \textit{«Неиспр. ЦУ Вn»}, \textit{«Ср. Пруж.
не заведены Вn»}, \textit{«Ср. Автомат ШП Вn»}.

\item
Также предусмотрена блокировка отключения от внешнего сигнала \textit{«Внеш. блок. откл. Вn»}, блокировка
включения от внешнего сигнала \textit{«Внеш. блок. вкл. Вn»} и блокировка управления (как включения, так и
отключения) от внешнего сигнала \textit{«Внеш. блок. упр. Вn»}.

\vspace{2mm}
\begin{figure}[H]
\centering
\includegraphics[width=0.9\textwidth,height=0.9\textheight,keepaspectratio]{\fbpath/761.3371 КСВ/img/BLK_OP_Q.pdf}
\captionof{figure}{Функциональный блок блокировки управления В1}\label{pict_ksv:BLK_OP_Q}
\end{figure}

\item
В устройстве предусмотрена программная фиксация команд или включенного положения выключателя
в зависимости от положения программного переключателя \textbf{«Реле фиксации Вn»}. В положении \textit{«РФК»}
фиксируется команда включения, в положении \textit{«РФП»} -- включенное положение. Квитирование сигнала
фиксации \textit{«РФ Вn»} может осуществляться от сброса сигнализации или подачей оперативной команды на
отключение. При отключении выключателя от защит и наличии сигнала \textit{«РФ Вn»} формируется сигнал о факте
аварийного отключения выключателя \textit{«Сигнал несоотв. Вn»}, необходимый для работы АПВ.

\vspace{2mm}
\begin{figure}[H]
\centering
\includegraphics[width=0.9\textwidth,height=0.9\textheight,keepaspectratio]{\fbpath/761.3371 КСВ/img/Relay_fix_Q.pdf}
\captionof{figure}{Функциональный блок реле фиксации В1}\label{pict_ksv:Relay_fix_Q}
\end{figure}


\small
\begin{longtable}{|>{\centering\arraybackslash}m{5.3cm}|>{\centering\arraybackslash}m{3.3cm}|>{\centering\arraybackslash}m{4.2cm}|>{\centering\arraybackslash}m{1.8cm}|>{\centering\arraybackslash}m{1cm}|}
\caption{Перечень уставочных параметров КСВ\hfill\vspace{-0.5\baselineskip}}\label{ksv:tbl1}\\ 
\hline
\rowcolor{gray!30}
Параметр & Условное обозначение на схеме & Значение/ Диапазон & Единица измерения & Шаг \\ 
\hline
\endfirsthead
\caption*{\hspace{3pt}\emph{Продолжение таблицы \ref{ksv:tbl1}\hfill\vspace{-0.5\baselineskip}}} \\ % сделано по ГОСТ 2.105 п.6.8.7
\hline
\rowcolor{gray!30}
Параметр & Условное обозначение на схеме & Значение/ Диапазон & Единица измерения & Шаг \\ 
%\hline
\endhead

%===m>КСВ|-
%===m
\end{longtable}
\normalsize

\end{enumerate}